% Created 2026-06-01 Mon 23:45
% Intended LaTeX compiler: pdflatex
\documentclass[11pt, letterpaper]{article}
\usepackage[utf8]{inputenc}
\usepackage[T1]{fontenc}
\usepackage{graphicx}
\usepackage{longtable}
\usepackage{wrapfig}
\usepackage{rotating}
\usepackage[normalem]{ulem}
\usepackage{amsmath}
\usepackage{amssymb}
\usepackage{capt-of}
\usepackage{hyperref}
\usepackage[margin=2.5cm]{geometry}
\usepackage[utf8]{inputenc}
\usepackage{palatino}
\usepackage{xcolor}
\author{Equipo Alpine White}
\date{\textit{{[}2026-05-04 Mon]}}
\title{Bitácora Semanal 13: Administración de Sistemas Unix/Linux}
\hypersetup{
 pdfauthor={Equipo Alpine White},
 pdftitle={Bitácora Semanal 13: Administración de Sistemas Unix/Linux},
 pdfkeywords={},
 pdfsubject={},
 pdfcreator={Emacs 30.2 (Org mode 9.7.39)}, 
 pdflang={English}}
\begin{document}

\maketitle
\setcounter{tocdepth}{2}
\tableofcontents

\section{Información General}
\label{sec:org76c6d5e}

\begin{itemize}
\item \textbf{Semana:} Semana 13
\item \textbf{Materia:} Administración de Sistemas Unix/Linux
\item \textbf{Docente:} Luis Enrique Serrano Gutiérrez
\item \textbf{Equipo:}

\begin{itemize}
\item Arreguín Salgado Gael Emiliano
\item López Pérez Mariana
\item Nieto Gallegos Isaac Julián
\end{itemize}
\end{itemize}

Durante esta semana se trabajó con tecnologías relacionadas con almacenamiento redundante, sistemas de archivos y compartición de recursos en red mediante NFS. La práctica permitió comprender cómo Linux administra arreglos RAID, integra múltiples sistemas de archivos y comparte almacenamiento entre equipos dentro de una red local.

El enfoque principal consistió en explorar configuraciones RAID utilizando LVM, comprender el funcionamiento de la paridad y reforzar conceptos relacionados con montaje de sistemas de archivos remotos mediante NFS.
\subsection{Responsabilidades}
\label{sec:org2207b10}

\begin{itemize}
\item Gael Emiliano: Investigación y documentación técnica relacionada con RAID, sistemas de archivos y configuraciones de almacenamiento redundate.

\item Mariana: Captura de temas vistos en laboratorio y ampliación de conceptos relacionados con NFS y montaje de recursos compartidos.

\item Isaac Julián: Revisión final, adaptación al formato org-mode y compilación del documento.
\end{itemize}

---
\section{RAID y almacenamiento redundante}
\label{sec:orgd983531}

\subsection{Conceptos nuevos}
\label{sec:orgbd832ee}

\subsubsection{¿Qué es RAID?}
\label{sec:org76aa902}

RAID (Redundant Array of Independent Disks) es una técnica que permite combinar varios discos duros o SSD para mejorar:

\begin{itemize}
\item Velocidad
\item Seguridad de datos
\item Tolerancia a fallos
\item Rendimiento general
\end{itemize}

RAID permite distribuir información entre varios dispositivos físicos según distintos niveles o estrategias.

---
\subsubsection{RAID 0}
\label{sec:org2a48b40}

RAID 0 divide los datos entre varios discos para incrementar considerablemente el rendimiento.

Ventajas:

\begin{itemize}
\item Mayor velocidad de lectura y escritura
\item Aprovechamiento completo del espacio
\end{itemize}

Desventajas:

\begin{itemize}
\item No existe redundancia
\item Si falla un disco, se pierde toda la información
\end{itemize}

Por esta razón RAID 0 se utiliza principalmente cuando el rendimiento es más importante que la seguridad de los datos.

---
\subsubsection{RAID 1}
\label{sec:org0d6311b}

RAID 1 utiliza una técnica de espejo (mirroring), duplicando la información entre dos discos.

Ventajas:

\begin{itemize}
\item Alta tolerancia a fallos
\item Protección de información
\item Recuperación rápida
\end{itemize}

Desventajas:

\begin{itemize}
\item Solo se aprovecha la mitad del espacio disponible
\end{itemize}

Este tipo de RAID es utilizado cuando la disponibilidad y la integridad de los datos son prioritarias.

---
\subsubsection{RAID 5}
\label{sec:org35d4cc4}

RAID 5 distribuye datos y paridad entre múltiples discos.

Características:

\begin{itemize}
\item Requiere mínimo 3 discos
\item Mantiene redundancia
\item Ofrece buen rendimiento
\item Optimiza espacio disponible
\end{itemize}

Gracias al uso de paridad, el sistema puede reconstruir información si uno de los discos falla.

---
\subsubsection{Paridad}
\label{sec:org1d97edb}

La paridad consiste en un cálculo matemático utilizado para reconstruir datos perdidos dentro de un arreglo RAID.

Generalmente se utiliza una operación XOR para generar información redundante distribuida entre discos.

Si un disco falla:

\begin{itemize}
\item Los datos faltantes pueden recalcularse
\item El sistema continúa funcionando
\item No se pierde información inmediatamente
\end{itemize}

---
\section{Sistemas de archivos}
\label{sec:orgcd55d52}

\subsection{Conceptos nuevos}
\label{sec:orga885e5d}

\subsubsection{ext4}
\label{sec:org26f669d}

ext4 es uno de los sistemas de archivos más utilizados en Linux debido a su estabilidad y rendimiento.

Características:

\begin{itemize}
\item Journaling
\item Buen rendimiento general
\item Alta compatibilidad
\item Bajo consumo de recursos
\end{itemize}

---
\subsubsection{XFS}
\label{sec:orgbeb2399}

XFS es un sistema de archivos avanzado optimizado para:

\begin{itemize}
\item Alto rendimiento
\item Grandes volúmenes de datos
\item Archivos de gran tamaño
\end{itemize}

Es ampliamente utilizado en servidores empresariales y sistemas de almacenamiento masivo.

---
\subsubsection{NTFS}
\label{sec:orgc475080}

NTFS es el sistema de archivos principal utilizado por Windows.

Características:

\begin{itemize}
\item Soporte para permisos
\item Journaling
\item Compresión
\item Cifrado
\item Manejo de archivos grandes
\end{itemize}

En Linux se requiere soporte adicional para escritura segura.

---
\subsubsection{ntfs-3g}
\label{sec:org21b3a89}

\texttt{ntfs-3g} es un controlador de código abierto que permite a Linux leer y escribir particiones NTFS.

Esto facilita:

\begin{itemize}
\item Compatibilidad con Windows
\item Transferencia de datos
\item Sistemas híbridos Linux/Windows
\end{itemize}

Durante la práctica fue necesario instalar manualmente los paquetes RPM de ntfs-3g debido a que RHEL no lo incluía por defecto.

---
\section{Compartición de recursos mediante NFS}
\label{sec:orgb223cbe}

\subsection{Conceptos de repaso}
\label{sec:org3f08d32}

\subsubsection{¿Qué es NFS?}
\label{sec:org507ef47}

NFS (Network File System) es un protocolo que permite compartir directorios y archivos entre sistemas Linux mediante red.

Gracias a NFS:

\begin{itemize}
\item Los recursos remotos pueden utilizarse como si fueran locales
\item Los archivos pueden centralizarse
\item Se facilita la administración en entornos multiusuario
\end{itemize}

---
\subsubsection{Servidor NFS}
\label{sec:org691ae16}

El servidor NFS es el equipo encargado de compartir directorios hacia otros sistemas.

La configuración de directorios compartidos se realiza mediante:

\begin{verbatim}
/etc/exports
\end{verbatim}

---
\subsubsection{Cliente NFS}
\label{sec:org45021aa}

El cliente NFS es el sistema que accede a los recursos compartidos proporcionados por el servidor.

---
\subsubsection{showmount -e}
\label{sec:org8e5f362}

El comando:

\begin{verbatim}
showmount -e servidor
\end{verbatim}

permite visualizar qué directorios está exportando un servidor NFS y verificar permisos de acceso.

---
\subsubsection{Montaje manual con mount}
\label{sec:org8b7769c}

Para conectar manualmente un recurso compartido NFS se utiliza:

\begin{verbatim}
sudo mount servidor:/ruta/compartida /mnt/nfs
\end{verbatim}

Requisitos:

\begin{itemize}
\item El servidor NFS debe estar activo
\item El recurso debe estar correctamente exportado
\item El cliente debe tener soporte NFS instalado
\item El directorio de montaje debe existir
\end{itemize}

---
\subsubsection{Verificación del montaje}
\label{sec:org69cc27c}

Después del montaje, se puede verificar el acceso mediante:

\begin{verbatim}
ls -la /mnt/nfs
\end{verbatim}

---
\section{LVM y administración flexible de almacenamiento}
\label{sec:org89dc2a2}

\subsection{Conceptos de repaso}
\label{sec:org9d94b42}

LVM continúa siendo una de las herramientas más importantes para abstraer almacenamiento físico y permitir administración dinámica de discos.

Componentes principales:

\begin{itemize}
\item PV (Physical Volume)
\item VG (Volume Group)
\item LV (Logical Volume)
\end{itemize}

Durante la práctica se reforzó el proceso de:

\begin{itemize}
\item Crear volúmenes físicos con \texttt{pvcreate}
\item Expandir grupos mediante \texttt{vgextend}
\item Aumentar capacidad usando \texttt{lvextend}
\item Redimensionar sistemas de archivos con \texttt{resize2fs}
\end{itemize}

---
\section{Herramientas complementarias}
\label{sec:orgd67b5cc}

\subsection{exportfs -av}
\label{sec:org00f2bf9}

Este comando aplica inmediatamente cambios realizados en \texttt{/etc/exports} sin necesidad de reiniciar completamente el servicio NFS.

\begin{verbatim}
exportfs -av
\end{verbatim}

---
\subsection{python3 -m http.server}
\label{sec:orgf8cf5c8}

Herramienta rápida utilizada para crear un servidor web temporal.

\begin{verbatim}
python3 -m http.server
\end{verbatim}

Durante la práctica fue utilizada para transferir paquetes RPM desde la máquina anfitriona hacia una máquina virtual aislada sin acceso a internet.

---
\section{Conceptos de repaso}
\label{sec:org32fcc7e}

\begin{itemize}
\item Linux utiliza abstracciones avanzadas para separar hardware físico y administración lógica.
\item RAID mejora disponibilidad y rendimiento dependiendo de la configuración utilizada.
\item NFS facilita el acceso compartido a recursos dentro de una red.
\item Los sistemas de archivos pueden coexistir sobre la misma infraestructura LVM sin conflictos.
\end{itemize}

---
\section{Máximas}
\label{sec:org1b850b4}

\begin{itemize}
\item “RAID 0 apuesta todo al rendimiento; si un disco falla, todo se pierde. La velocidad tiene un precio.”

\item “RAID 5 es el equilibrio que busca todo administrador: rendimiento, redundancia y espacio razonable en un solo arreglo.”

\item “Un Volume Group no es un disco; es una reserva de espacio que puede crecer, encogerse y reorganizarse sin que el sistema lo note.”

\item “ext4, xfs y NTFS sobre el mismo servidor no es redundancia; es demostrar que LVM no le importa qué sistema de archivos vive encima.”
\item “pvcreate convierte en materia prima a un disco para que LVM pueda moldearlo a su voluntad.”
\end{itemize}

---
\section{Sección libre}
\label{sec:orgace2eff}

\begin{itemize}
\item Los arreglos RAID permiten adaptar el almacenamiento según necesidades específicas de rendimiento o seguridad.

\item \texttt{showmount -e} es una herramienta esencial para verificar exportaciones NFS activas.
\end{itemize}

\begin{verbatim}
showmount -e
\end{verbatim}

\begin{itemize}
\item \texttt{ntfs-3g} mejora considerablemente la interoperabilidad entre Linux y Windows.
\item NFS permite compartir recursos de red de manera transparente para usuarios y aplicaciones.

\item LVM facilita el crecimiento dinámico del almacenamiento sin necesidad de detener servicios.
\end{itemize}

---
\section{Realimentación y reflexión de aprendizajes}
\label{sec:org5a5319f}

Durante esta semana se comprendió cómo Linux integra múltiples tecnologías de almacenamiento y compartición de recursos dentro de un mismo entorno administrativo.

La práctica con RAID permitió observar cómo diferentes estrategias priorizan rendimiento, redundancia o equilibrio entre ambos. RAID 0 mostró la importancia del rendimiento, RAID 1 reforzó conceptos de tolerancia a fallos y RAID 5 introdujo el uso práctico de paridad para recuperación de datos.

El trabajo con NFS permitió entender cómo Linux comparte recursos remotos de forma transparente mediante montaje de directorios, facilitando administración centralizada y colaboración entre sistemas.

Finalmente, el uso combinado de LVM, múltiples sistemas de archivos y herramientas como ntfs-3g mostró cómo Linux puede integrar distintos tipos de almacenamiento y plataformas sin perder flexibilidad ni estabilidad.
\end{document}
